\section*{When averaging graphs creates a mode that no context has}
\addcontentsline{toc}{section}{When averaging graphs creates a mode that no context has}

Consider three positions and two equally likely sequence contexts.  In context $A$, only positions
1 and 2 are connected; in context $B$, only positions 2 and 3 are connected.  Their unnormalized
Laplacians are
\begin{equation}
L_A=
\begin{pmatrix}
1&-1&0\\
-1&1&0\\
0&0&0
\end{pmatrix},
\qquad
L_B=
\begin{pmatrix}
0&0&0\\
0&1&-1\\
0&-1&1
\end{pmatrix}.
\label{eq:book-context-laplacians}
\end{equation}
Each graph consists of one edge and one isolated vertex, so both spectra are
\begin{equation}
\operatorname{spec}(L_A)
=\operatorname{spec}(L_B)
=\{0,0,2\}.
\end{equation}

The mean Laplacian is
\begin{equation}
\overline L
=\frac12(L_A+L_B)
=\begin{pmatrix}
1/2&-1/2&0\\
-1/2&1&-1/2\\
0&-1/2&1/2
\end{pmatrix},
\end{equation}
with spectrum
\begin{equation}
\operatorname{spec}(\overline L)=\{0,1/2,3/2\}.
\label{eq:book-mean-laplacian-spectrum}
\end{equation}
The mean graph is connected even though no context contains both edges simultaneously.  Averaging
has created a path $1\!-\!2\!-\!3$ that exists only across contexts.  Correspondingly, the second
zero mode present in every individual context disappears from the mean.

This tiny example separates three claims.  The mean context spectrum is $\{0,0,2\}$.  The spectrum
of the mean graph is $\{0,1/2,3/2\}$.  The distribution of zero-mode projectors records which vertex
is isolated in each context.  None determines the others.

For protein responses, the distinction can be biological.  Two alternative sequence backgrounds
may support mutually exclusive communication paths.  Their mean graph suggests a connected route
through all three positions, while no individual protein realizes that route.  A spectral analysis
of $\E_C A(C)$ therefore answers a population-level aggregation question, not a statement about a
typical sequence.  Context-specific projectors or their distribution must be retained when that
difference matters.
